\section{Boundary rigidity for the canonical atom spectra}

Let \(F\) be a connected reduced obligatory triple system with \(m\)
hyperedges and \(n\) points, and put
\[
 s=n-m.
\]
Define
\[
 \alpha(s)=
 \begin{cases}
 s,& s\text{ even},\\
 s+1,& s\text{ odd},
 \end{cases}
 \qquad
 \delta(s)=
 \left\lfloor\frac{(s-1)^2}{4}\right\rfloor+1.
\]
For positive integers \(p,q,r\), let \(\Theta(p,q,r)\) denote the graph
formed by three internally disjoint paths of lengths \(p,q,r\) with common
endpoints.

\begin{theorem}[Atomic lower-bound rigidity]
Let \(F\) be one-point indecomposable and not a single triple.  Then
\(s\ge4\) and \(m\ge\alpha(s)\).  Equality is classified as follows.
\begin{enumerate}
\item If \(s\) is even, then \(m=s\) if and only if
\[
 F\cong C_s^+.
\]
\item If \(s\) is odd, then \(m=s+1\) if and only if
\[
 F\cong \Theta(2a,2b,2c)^+
\]
for positive integers \(a\le b\le c\) with
\[
 a+b+c=\frac{s+1}{2}.
\]
\end{enumerate}
Hence, for odd \(s\ge5\), the number of isomorphism types at the atomic lower
boundary is
\[
 \left\lfloor
 \frac{\bigl((s+1)/2\bigr)^2+3}{12}
 \right\rfloor.
\]
\end{theorem}

\begin{proof}
By the canonical atom theorem, \(F\cong J^+\) for a finite 2-connected simple
bipartite graph \(J\) with \(|V(J)|=s\) and \(|E(J)|=m\).  Every vertex of
\(J\) has degree at least two.

If \(s\) is even and \(m=s\), every vertex has degree two, so \(J\) is the
cycle \(C_s\).

If \(s\) is odd and \(m=s+1\), then
\[
 \sum_{v\in V(J)}(\deg v-2)=2.
\]
The degree-excess pattern consisting of one degree-four vertex and all other
vertices of degree two would make that degree-four vertex a cut vertex.
Therefore \(J\) has exactly two degree-three vertices and every other vertex
has degree two.  Suppressing the degree-two chains gives three internally
disjoint paths between the degree-three vertices.  Bipartiteness forces their
lengths to have the same parity, and their sum \(m=s+1\) is even, so all three
lengths are even.  This gives \(\Theta(2a,2b,2c)\), and the converse is
immediate.

The isomorphism type is determined by the unordered triple
\(\{a,b,c\}\).  The number of positive triples \(a\le b\le c\) with sum \(N\)
is
\[
 \left\lfloor\frac{N^2+3}{12}\right\rfloor.
\]
Take \(N=(s+1)/2\).
\end{proof}

\begin{theorem}[Atomic upper-bound rigidity]
Let \(F\) be as above.  Then
\[
 m\le \left\lfloor\frac{s^2}{4}\right\rfloor,
\]
with equality if and only if
\[
 F\cong K_{\lfloor s/2\rfloor,\lceil s/2\rceil}^{+}.
\]
\end{theorem}

\begin{proof}
If the bipartition sizes of the 2-connected core are \(a+b=s\), then
\[
 m\le ab\le\left\lfloor\frac{s^2}{4}\right\rfloor.
\]
Equality forces the balanced complete bipartite graph.
\end{proof}

\begin{lemma}[Cut-vertex extremisers]
Let \(J\) be a connected simple bipartite graph on \(s\ge3\) vertices with a
cut vertex.  Then
\[
 |E(J)|\le
 \left\lfloor\frac{(s-1)^2}{4}\right\rfloor+1.
\]
Equality holds if and only if \(J\) is obtained from the balanced complete
bipartite graph on \(s-1\) vertices by adjoining one leaf at an arbitrary
vertex.
\end{lemma}

\begin{proof}
Let the bipartition be \(A\sqcup B\), with \(|A|=a\), \(|B|=b\), and suppose
a cut vertex \(x\) lies in \(A\).  If the components of \(J-x\) have
bipartition sizes \((a_i,b_i)\), then every component contains a vertex of
\(B\), and there are at least two components.  Hence
\[
 |E(J)|
 \le b+\sum_i a_i b_i
 \le b+(a-1)(b-1)
 =ab-a+1.
\]
The symmetric estimate holds when the cut vertex lies in \(B\).  Thus
\[
 |E(J)|\le ab-\min(a,b)+1
 \le \left\lfloor\frac{(s-1)^2}{4}\right\rfloor+1.
\]
Equality forces one opposite-side singleton component and a complete
bipartite large component, except for the star cases \(s=3,4\), which have the
same description after deleting one leaf.  The complete bipartite remainder
must be balanced.  The converse is immediate.
\end{proof}

\begin{theorem}[Decomposable upper-bound rigidity]
Let \(F\) be one-point decomposable.  Then
\[
 m\le\delta(s).
\]
Equality holds if and only if there exist \(a,b\ge1\) such that
\[
 a+b=s-1,\qquad |a-b|\le1,
\]
and \(F\) is a one-point amalgamation of \(K_{a,b}^+\) and one additional
triple.  The amalgamation point may be any point of either system.
\end{theorem}

\begin{proof}
The canonical atom forest gives a connected bipartite parameter shadow with a
cut vertex.  The preceding lemma gives the bound and identifies the equality
shadow as a balanced complete bipartite graph plus one leaf.  Translating its
blocks through the canonical atom theorem gives \(K_{a,b}^+\) and one
single-triple piece.  Conversely, such an amalgamation has
\[
 m=ab+1,\qquad s=a+b+1,
\]
and therefore attains \(\delta(s)\).
\end{proof}

\begin{corollary}[Boundary table]
Excluding the single triple, the extremal connected boundaries are:
\[
\begin{array}{c|c|c}
\text{regime} & \text{edge count} & \text{structure}\\ \hline
\text{decomposable lower} & m=s-1
  & \text{all atoms are single triples}\\
\text{atomic lower, }s\text{ even} & m=s
  & C_s^+\\
\text{atomic lower, }s\text{ odd} & m=s+1
  & \Theta(2a,2b,2c)^+\\
\text{decomposable upper} & m=\delta(s)
  & K_{a,b}^+\text{ plus one triple}\\
\text{atomic upper} & m=\lfloor s^2/4\rfloor
  & K_{\lfloor s/2\rfloor,\lceil s/2\rceil}^+
\end{array}
\]
where \(a+b=s-1\) and \(|a-b|\le1\) in the decomposable upper row.
\end{corollary}
